\section{서 론}
\label{sec:introduction}

양자컴퓨팅 기술의 발전으로 기존 공개키 암호체계의 안전성에 대한 우려가
증가함에 따라, 격자 문제에 기반한 양자내성암호 연구가 활발히 진행되고
있다. 미국 국립표준기술연구소(NIST)는 격자 기반 전자서명
\code{CRYSTALS-Dilithium}을 ML-DSA로 표준화하였으며
~\cite{dilithium,fips204,nist_pqc}, 국내 KpqC 공모에서는 격자 기반
전자서명 \code{HAETAE}가 최종 후보로 선정되었다. \code{HAETAE}는
Module-LWE 및 Module-SIS 문제에 안전성을 기반하는 Fiat--Shamir with
Aborts 방식의 전자서명으로, 초구 샘플링과 bimodal 거부 샘플링을 이용해
서명 크기와 구현 효율을 개선하였다~\cite{haetae}.

그러나 암호 알고리즘의 수학적 안전성이 실제 구현 환경의 안전성을 직접
보장하지는 않는다. 오류주입 공격(Fault Injection Attack)은 클럭, 전압,
전자기파 또는 레이저 등을 이용하여 암호 장치의 중간값이나 프로그램
흐름을 변조하는 공격이다~\cite{yang2020,oflynn2016,skorobogatov2002,
dumont2019}. 공격자는 오류가 포함된 출력과 정상 출력을 분석하여 비밀정보를
복구하거나 보안 검사를 우회할 수 있다. Boneh 등은 RSA 서명 과정에 발생한
오류로부터 비밀키를 복구할 수 있음을 보였으며~\cite{boneh1997}, 이후
오류주입 공격의 대상은 대칭키 암호와 공개키 암호 전반으로 확대되었다
~\cite{barel2006,ravi2022survey}.

격자 기반 전자서명에서도 거부 샘플링 우회, 결정론적 난수 재사용, 명령어
생략 및 NTT 연산 변조를 이용한 공격이 보고되었다~\cite{espitau2016,
bruinderink2018,ravi2019determinism,ravi2023ntt}. 특히 서명 응답이
일회용 난수와 비밀키의 선형 결합으로 구성되는 경우, 덧셈 항의 일부가
제거되면 정상적으로는 가려져야 할 비밀정보의 관계가 출력에 남을 수 있다.
ML-DSA에 대해서도 단일 또는 소수의 오류를 이용한 비밀정보 복원 가능성이
연구되었다~\cite{krahmer2024,elghamrawy2023,jendral2024}.

\code{HAETAE}의 서명 응답은 다음과 같이 계산된다.
\begin{equation}
  \zvec=\yvec+(-1)^b c\svec_1 ,
  \label{eq:intro-response}
\end{equation}
여기서 $\yvec$는 일회용 난수, $b$는 bimodal 부호, $c$는 공개 챌린지,
$\svec_1$은 서명 비밀 성분이다. 식~\eqref{eq:intro-response}에서
$c\svec_1$ 또는 $\yvec$가 제거되면 정상 응답과 다른 대수적 관계가
형성된다. 이러한 관계는 소프트웨어 수준에서는 간단한 항 제거로 표현할 수
있지만, 실제 물리 글리치가 동일한 효과를 생성하는지는 컴파일된 연산 구조와
버퍼 상태에 따라 달라질 수 있다.

기존 오류주입 연구는 대체로 특정 오류가 발생했을 때의 키 복구 가능성이나
대응기법의 탐지 성능에 초점을 두었다. 반면 소프트웨어에서 정의한 오류
효과가 실제 구현에서 어떤 중간값으로 나타나며, 그 값이 직렬화된 서명까지
전파되는지를 단계적으로 구분한 연구는 상대적으로 부족하다. 특히 의미론적
오류 에뮬레이션의 성공을 곧바로 물리 공격 성공으로 해석하면 공격 가능성을
과대평가할 수 있다. 반대로 특정 글리치 실험에서 목표 오류가 관측되지
않았다는 이유만으로 대수적 취약성이 제거되었다고 판단하는 것 역시
적절하지 않다.

따라서 본 논문에서는 소프트웨어 오류 에뮬레이션으로 \code{HAETAE}
전체 서명 경로의 취약 지점을 체계화한다. 이 가운데 T1에 대해서는
STM32F405 클럭 글리치 실험을 수행하여 물리적 블록 제거와 비밀 성분
복원을 검증한다. T2의 미수정 융합 구현에 대한 미관측 결과는 대수적
취약성과 물리적 실현 가능성의 차이를 설명하는 비교 자료로 사용한다.

\subsection*{Contributions}
\begin{itemize}[leftmargin=1.2em,itemsep=0.35em,topsep=0.2em]
  \item \code{HAETAE}의 전체 서명 경로에서 발생 가능한 오류 지점을
        상류 연산, 응답 연산 및 거부 판정으로 구분하고, 오류 효과와
        공격에 필요한 출력 및 검증 범위를 체계화하였다. 상류 지점과
        T1, T2 및 RB의 의미론적 오류 효과는 전체 서명 코드의 SW
        에뮬레이션으로 확인하였다.
  \item 응답에서 비밀키 곱 항이 제거되는 T1과 일회용 난수 항이 제거되는
        T2의 대수적 복원 관계를 도출하였다. T2는 단일 오류 응답에서
        $\svec_1$을 역산할 수 있는 반면, T1은 동일 nonce의 정상 응답과
        결함 응답 사이의 차분이 필요함을 보였다. T2 SW 실험에서는
        25/30회 오류 응답이 방출되었고, 각 경우에 $\svec_1$의
        768개 계수가 모두 일치하였다.
  \item STM32F405의 T1 인과대조 구현에서 실제 클럭 글리치로 다항식
        블록 제거를 발생시켰다. 공격 구간을 식별한 후 두 개의 유용한
        결함 응답을 누적하여 $\svec_1$ 768/768 계수를 복원하였고,
        별도 실행에서도 594회 이내에 이를 재현하였다.
  \item 미수정 융합 응답 구현에 650회의 클럭 글리치를 주입하였으나
        깨끗한 T2는 관측되지 않았다. T1 물리 결과와 이를 비교하여,
        버퍼 초기 상태와 반복 구조가 오류 효과의 물리적 실현에 영향을
        줄 수 있음을 실험적으로 시사한다.
\end{itemize}

본 논문의 구성은 다음과 같다. 2장에서는 \code{HAETAE}의 서명 구조와
관련 오류주입 연구를 설명한다. 3장에서는 오류모델과 응답 연산의 취약성을
분석한다. 4장에서는 소프트웨어 및 물리 실험 결과와 그 의미 및 한계를
논의한다. 마지막으로 5장에서 결론을 맺는다.
